\documentclass{article}
\usepackage{amsfonts,amsthm,amsmath,amssymb,mathtools}
\usepackage{mathrsfs}
\usepackage{enumitem}
\usepackage{tikz-cd}

\setcounter{section}{0}
\renewcommand{\thesection}{\S\arabic{section}}
\renewcommand{\thesubsection}{\arabic{section}}

\newtheorem{thm}{Theorem}
\newenvironment{theorem}[1]{
	\renewcommand{\thethm}{#1}
	\begin{thm}
		}{\end{thm}}

\newtheorem{lma}{Lemma}
\newenvironment{lemma}[1]{
	\renewcommand{\thelma}{#1}
	\begin{lma}
		}{\end{lma}}

\newcounter{section-number}
\setcounter{section-number}{3}

\theoremstyle{definition}
\newtheorem{ex}{}[section-number]
\newenvironment{exercise}[1]{
  \renewcommand{\theex}{\arabic{section-number}.\arabic{ex}}
	\begin{ex}
		}{\end{ex}}

\title{Exercises I.3}
\author{Connor Mooneyhan}
\date{June 12, 2024}

\begin{document}

\maketitle

\begin{ex}
	Let $\mathsf{C}$ be a category.
	Consider a structure $\mathsf{C}^{op}$ with \begin{itemize}
		\item $\mathrm{Obj}(\mathsf{C}^{op}) \coloneq \mathrm{Obj}(\mathsf{C})$;
		\item for $A, B$ objects of $\mathsf{C}^{op}$ (hence objects of $\mathsf{C}$), $\mathrm{Hom}_{\mathsf{C}^{op}}(A, B) \coloneq \mathrm{Hom}_{\mathsf{C}}(B, A)$.
	\end{itemize}
	Show how to make this into a category (that is, define composition of morphisms in $\mathsf{C}^{op}$ and verify the properties listed in \S 3.1).

	Intuitively, the 'opposite' category $\mathsf{C}^{op}$ is simply obtained by 'reversing all the arrows' in $\mathsf{C}$.
\end{ex}
\begin{proof}
	Let $A, B, C$ be objects of $\mathsf{C}^{op}$ and $f \in \mathrm{Hom}_{\mathsf{C}^{op}}(A, B)$, $g \in \mathrm{Hom}_{\mathsf{C}^{op}}(B, C)$.
	Introduce notation $\phi'$ for any morphism $\phi$ in $\mathsf{C}^{op}$ to denote its analogue in $\mathsf{C}$.
	Then we define the composition $gf$ as \begin{equation*}
		gf \coloneq f'g' \in \mathrm{Hom}_{\mathsf{C}}(C, A).
	\end{equation*}
	Thus our notation indicates that $(gf)' = f'g'$.

	Identity morphisms are preserved, of course, since \begin{equation*}
		f\mathrm{id}_A \coloneq \mathrm{id}_A'f' = f' \eqcolon f,
	\end{equation*}
	and similarly for $\mathrm{id}_A f$.
	Composition is associative since given one more object $D$ in $\mathsf{C}^{op}$ and morphism $h \in \mathrm{Hom}_{C^{op}}(C, D)$, we have \begin{equation*}
		(hg)f \eqcolon f'(hg)' = f'(g'h') = (f'g')h' = (gf)'h' \coloneq h(gf).
	\end{equation*}
\end{proof}

\begin{ex}
	If $A$ is a finite set, how large is $\mathrm{End}_{\mathsf{Set}}(A)$?
\end{ex}
\begin{proof}
	Clearly $\mathrm{End}_{\mathsf{Set}}(A) = A^A$, so \begin{equation*}
		|\mathrm{End}_{\mathsf{Set}}(A)| = |A^A| = |A|^{|A|}
	\end{equation*}
\end{proof}

\begin{ex}
	Formulate precisely what it means to say that $1_a$ is an identity with respect to composition in Example 3.3, and prove this assertion.
\end{ex}
\begin{proof}
	Given $f \in \mathrm{Hom}(a, b)$, to say that $1_a$ and $1_b$ are identities with respect to composition is to say that they preserve the source and target when composed with $f$.
	This is the only way they could produce a different function, since $f$ is the only member of $\mathrm{Hom}(a, b)$.

	Of course, $f 1_a, 1_b f \in \mathrm{Hom}(a, b)$ by definition, so this is true.
\end{proof}

\begin{ex}
	Can we define a category in the style of Example 3.3 using the relation $<$ on the set $\mathbb{Z}$?
\end{ex}
\begin{proof}
	No.
	Since $<$ is not reflexive, $\mathrm{Hom}(A, A)$ would be empty for all objects $A$, so there is nothing to define as the identity $1_A$.
\end{proof}

\begin{ex}
	Explain in what sense Example 3.4 is an instance of the categories considered in Example 3.3.
\end{ex}
\begin{proof}
	The subset relation $\subseteq$ is both reflexive and transitive, and thus satisfies the requirements of the relation $\sim$ in Example 3.3.
\end{proof}

\begin{ex}
	(Assuming some familiarity with linear algebra.)
	Define a category $\mathsf{V}$ by taking $\mathrm{Obj}(\mathsf{V}) = \mathbb{N}$ and letting $\mathrm{Hom}_{\mathsf{V}}(n, m) =$ the set of $m \times n$ matrices with real entries, for all $n, m \in \mathbb{N}$.
	(I will leave the reader the task of making sense of a matrix with 0 rows or columns.)
	Use products of matrices to define composition.
	Does this category 'feel' familiar?
\end{ex}
\begin{proof}
	We define $\mathsf{V}$ by \begin{itemize}
		\item $\mathrm{Obj}(\mathsf{V}) = \mathbb{N}$
		\item $\mathrm{Hom}_{\mathsf{V}}(n, m) =$ the set of $m \times n$ matrices with real entries.
	\end{itemize}
	This definition is unclear in the case of a matrix with either 0 rows, 0 columns, or both.
	Composition is defined by matrix multiplication.
	Given $A$ an $n_2 \times n_1$ matrix and $B$ an $n_3 \times n_2$ matrix for some $n_1, n_2, n_3 \in \mathbb{N}$, we define $BA = E$ if any of $n_1, n_2,$ or $n_3$ are 0, where $E$ represents the $0 \times 0$ matrix.

	Composition is associative because matrix multiplication is associative, and we define the identity morphism $1_n$ to be the diagonal $n \times n$ matrix with $1$ in each place along the diagonal.
	It is a basic result of linear algebra that this acts as an identity with respect to matrix multiplication, and thus with respect to our composition in $\mathsf{V}$.

	This category does not remind me of anything.
	Sorry Paolo.
\end{proof}

\begin{ex}
	Define carefully objects and morphisms in Example 3.7, and draw the diagram corresponding to composition.
\end{ex}
\begin{proof}
	Given a category $\mathsf{C}$, we can notate the coslice category corresponding to it with respect to a fixed object $A$ by $\mathsf{C}^A$.
	The objects in $\mathsf{C}^A$ are exactly the morphisms in $\mathsf{C}$ with source $Z$, where $Z$ is any object of $\mathsf{C}$, and target $A$.
	Thus, an object $f$ in $\mathsf{C}^A$ corresponds to the diagram \begin{equation*}
		\begin{tikzcd}
			Z \\
			A \arrow[u, "f"]
		\end{tikzcd}
	\end{equation*}
	in $\mathsf{C}$.
	A morphism in $\mathsf{C}^A$ from objects $f_1: A \to Z_1$ to $f_2: A \to Z_2$ is defined to be a morphism $\phi$ in $\mathsf{C}$ from $Z_1$ to $Z_2$ such that $\phi f_1 = f_2$.
	In other words, it's a commutative diagram \begin{equation*}
		\begin{tikzcd}
			Z_1 \arrow[rr, "\phi"] & & Z_2 \\
			& A \arrow[ul, "f_1"] \arrow[ur, "f_2"'] &
		\end{tikzcd}
	\end{equation*}
	The commutative diagram corresponding to composition, then, looks like this when provided another object $f_3: A \to Z_3$: \begin{equation*}
		\begin{tikzcd}
			Z_1 \arrow[r, "\phi_1"] & Z_2 \arrow[r, "\phi_2"] & Z_3 \\
			& A \arrow[ul, "f_1"] \arrow[u, "f_2"] \arrow[ur, "f_3"']
		\end{tikzcd}
	\end{equation*}
\end{proof}

\begin{ex}
	A \textit{subcategory} $\mathsf{C}'$ of a category $\mathsf{C}$ consists of a collection of objects of $\mathsf{C}$ with sets of morphisms $\mathrm{Hom}_{\mathsf{C}'}(A, B) \subseteq \mathrm{Hom}_{\mathsf{C}}(A, B)$ for all objects $A, B$ in $\mathrm{Obj}(\mathsf{C}')$, such that identities and compositions in $\mathsf{C}$ make $\mathsf{C}'$ into a category.
	A subcategory $\mathsf{C}'$ is \textit{full} if $\mathrm{Hom}_{\mathsf{C}'}(A, B) = \mathrm{Hom}_{\mathsf{C}}(A, B)$ for all $A, B$ in $\mathrm{Obj}(\mathsf{C}')$.
	Construct a category of \textit{infinite sets} and explain how it may be viewed as a full subcategory of $\mathsf{Set}$.
\end{ex}
\begin{proof}
	Define the category $\mathsf{Set_\infty}$ by making the objects the class of infinite sets and the morphisms functions between infinite sets.
	Clearly, each object is a set, and thus is an object in $\mathsf{Set}$, and each morphism $f \in \mathrm{Hom}_\mathsf{Set_\infty}(A, B)$ is a function from sets $A$ to $B$, and thus is in $\mathrm{Hom}_{\mathsf{Set}}(A, B)$.
	Also, identities are the same and composition is preserved, so it is indeed a subcategory.
	To show that it is full, we take any morphism $f \in \mathrm{Hom}_{\mathsf{Set}}(A, B)$ where $A, B$ are in $\mathrm{Obj}_{\mathsf{Set_\infty}}$, and clearly $f \in \mathrm{Hom}_{\mathsf{Set}_\infty}$.
\end{proof}

\begin{ex}
	An alternative to the notion of \textit{multiset} introduced in \S 2.2 is obtained by considering sets endowed with equivalence relations; equivalent elements are taken to be multiple instances of elements 'of the same kind'.
	Define a notion of morphism between such enhanced sets, obtaining a category $\mathsf{MSet}$ containing (a 'copy' of) $\mathsf{Set}$ as a full subcategory.
	(There may be more than one reasonable way to do this! This is intentionally an open-ended exercise.)
	Which objects in $\mathsf{MSet}$ determine ordinary multisets as defined in \S 2.2 and how?
	Spell out what a morphism of multisets would be from this point of view.
	(There are several natural notions of morphisms of multisets. Try to define morphisms in $\mathsf{MSet}$ so that the notion you obtain for ordinary multisets captures your intuitive understanding of these objects.)
\end{ex}
\begin{proof}
	Define the objects of $\mathsf{MSet}$ to be $S/{\sim_S}$ where $S$ is a set and $\sim_S$ is an equivalence relation on that set.
	Morphisms shall be defined as simple functions between multisets.

	It is clear that for sets $S_1$ and $S_2$, the $=$ relation determines objects $S_1/{=}$ and $S_2/{=}$ in $\mathsf{MSet}$ such that any function $f: S_1 \to S_2$ corresponds to $\phi: S_1/{=} \to S_2/{=}$ where $\phi(\left\lbrace a \right\rbrace) = \left\lbrace f(a) \right\rbrace$.

	Taking on the notation of \S 2.2, you would need $|[a]_\sim| = m(a)$ for each $a \in A$, given $\sim$ some equivalence relation.

	Note that morphisms of multisets under this view would need to send 'each copy' of an element to the same place.
	This may or may not be ideal, but I really have no context for why I would need multisets, so it's unclear whether this should fit with applications or not.
\end{proof}

\begin{ex}
	Since the objects of a category $\mathsf{C}$ are not (necessarily interpreted as) sets, it is not clear how to make sense of a notion of 'subobject' in general.
	In some situations it \textit{does} make sense to talk about subobjects, and the subobjects of any given object $A$ in $\mathsf{C}$ are in one-to-one correspondence with the morphisms $A \to \Omega$ for a fixed, special object $\Omega$ of $\mathsf{C}$, called a \textit{subobject classifier}.
	Show that $\mathsf{Set}$ has a subobject classifier.
\end{ex}
\begin{proof}
	'The' subobject classifier in $\mathsf{Set}$ is 'the' two-element set (up to isomorphism), which can be thought of as $\left\lbrace 0, 1 \right\rbrace$.
	If we denote $2^A$ to be all the functions from a set $A$ to $\left\lbrace 0, 1 \right\rbrace$, then we have already shown that $2^A$ is in one-to-one correspondence with the subsets of $A$ in Exercise 2.11.
\end{proof}

\begin{ex}
	Draw the relevant diagrams and define composition and identities for the category $\mathsf{C}^{A, B}$ mentioned in Example 3.9.
	Do the same for the category $\mathsf{C}^{\alpha, \beta}$ mentioned in Example 3.10.
\end{ex}
\begin{proof}
	The objects of $\mathsf{C}^{A, B}$ are diagrams \begin{equation*}
		\begin{tikzcd}
			A \arrow[dr, "f"] \\
			& Z \\
			B \arrow[ur, "g"']
		\end{tikzcd}
	\end{equation*}
	where $Z$ is any object of $\mathsf{C}$, and the morphisms \begin{equation*}
		\begin{tikzcd}
			A \arrow[dr, "f_1"]  & \\
			& Z_1 \\
			B \arrow[ur, "g_1"'] &
		\end{tikzcd}
		\longrightarrow
		\begin{tikzcd}
			A \arrow[dr, "f_2"]  & \\
			& Z_2 \\
			B \arrow[ur, "g_2"'] &
		\end{tikzcd}
	\end{equation*}
	are commutative diagrams \begin{equation*}
		\begin{tikzcd}
			A \arrow[dr, "f_1"'] \arrow[drr, bend left, "f_2"] \\
			& Z_1 \arrow[r, "\sigma"] & Z_2 \\
			B \arrow[ur, "g_1"] \arrow[urr, bend right, "g_2"'] \\
		\end{tikzcd}
	\end{equation*}
	The identity morphism is found by simply taking $\sigma$ to be the identity in $\mathsf{C}$.

	Moving to $\mathsf{C}^{\alpha, \beta}$ where $\alpha: C \to A$, $\beta: C \to B$, then the objects are commutative diagrams \begin{equation*}
		\begin{tikzcd}
      & A \arrow[dr, "f"] \\
      C \arrow[ur, "\alpha"] \arrow[dr, "\beta"'] & & Z \\
      & B \arrow[ur, "g"'] \\
		\end{tikzcd}
	\end{equation*}
  and the morphisms are commutative diagrams \begin{equation*}
    \begin{tikzcd}
      & A \arrow[dr, "f_1"] \arrow[drr, bend left, "f_2"] \\
      C \arrow[ur, "\alpha"] \arrow[dr, "\beta"'] & & Z_1 \arrow[r, "\sigma"] & Z_2 \\
      & B \arrow[ur, "g_1"'] \arrow[urr, bend right, "g_2"']\\
    \end{tikzcd}
  \end{equation*}
  Identities are once again found by making $\sigma$ the identity morphism from $\mathsf{C}$.
\end{proof}

\end{document}
